\section{Introduction}
\label{sec:intro}

Modern language models are deployed in conversations that routinely span dozens of turns, yet the alignment techniques that make them safe and useful---supervised fine-tuning (SFT), reinforcement learning from human feedback (RLHF), and direct preference optimization (DPO)---are trained and evaluated predominantly on single-turn or short multi-turn interactions. A growing body of evidence suggests this mismatch matters: multi-turn performance degrades by an average of 39\% compared to single-turn equivalents~\citep{laban2025lost}, instruction retention plummets even in frontier models~\citep{sirdeshmukh2025multichallenge}, and safety alignment can be systematically eroded through sequential interaction~\citep{li2026multiturn_jailbreak}. But \emph{why} does this degradation occur?

\para{Our main question.} We hypothesize that multi-turn degradation is, at least partly, a consequence of \emph{regression to the prior}: as conversation context accumulates, the aligned model's output distribution drifts back toward the base model's pre-training distribution. This hypothesis is motivated by three observations from prior work. First, alignment is superficial---only 5--8\% of output tokens are meaningfully affected by alignment tuning, and these are predominantly stylistic~\citep{lin2024urial}. Second, safety alignment localizes to roughly 1.4\% of model neurons, and these neurons are fragile under distribution shift~\citep{li2024ssah}. Third, within individual responses, the KL divergence between aligned and base models decreases monotonically over token position~\citep{lin2024urial}, suggesting that the alignment signal fades as generation progresses. If this fading extends across turns, we would expect aligned model behavior to become progressively more base-like over the course of a conversation.

\para{Gap in existing work.} Despite extensive research on both the superficiality of alignment and the fragility of multi-turn performance, no prior work has directly measured whether aligned model outputs \emph{converge toward base model distributions} as turn count increases. Studies of multi-turn degradation focus on task accuracy~\citep{laban2025lost, myung2026conversational} or attack success rates~\citep{li2026multiturn_jailbreak}, while studies of alignment depth analyze single-turn or single-response phenomena~\citep{lin2024urial, li2024ssah, vergara2026superficial}. The critical missing link is a direct measurement of instruct--base similarity across conversation depth.

\para{Our approach.} We address this gap through two complementary experiments. In Experiment~1, we compare next-token probability distributions between \qwenbase (base) and \qweninst (instruct) given identical conversation prefixes of varying depth (0--12 turns), measuring KL divergence, Jensen-Shannon divergence, and top-$k$ token overlap. A \concat control, which presents the same conversational content without turn structure, isolates the effect of multi-turn formatting from information accumulation. In Experiment~2, we deploy behavioral alignment probes---safety refusal, formatting compliance, and instruction retention---at different turn positions (0--24) in conversations with \gptmini and \gptnano via the OpenAI API.

\para{Preview of results.} We find partial support for the regression-to-prior hypothesis with important nuances. Superficial alignment behaviors degrade significantly: uppercase formatting compliance drops from 100\% to 25\% ($p<0.001$, $d=2.45$) and word limit compliance drops from 100\% to 81\% ($p<0.001$, $d=1.10$) over 24 turns. Safety alignment, however, remains robust through all 24 turns tested. Token-level analysis reveals a two-phase pattern---alignment signal strengthens in early turns (0$\to$2) as chat formatting activates, then gradually weakens through turns 2--12---suggesting that regression to the prior competes with a format-activation effect. Unexpectedly, multi-turn conversation structure produces systematically higher instruct--base divergence than concatenated equivalents ($\Delta \approx 3.1$ nats, $p<0.001$ at all turns), indicating that turn structure itself interacts with alignment mechanisms.

We make the following contributions:
\begin{itemize}[leftmargin=*,itemsep=0pt,topsep=2pt]
    \item We conduct the first direct measurement of instruct--base distribution similarity across conversation turns, revealing a two-phase activation-then-regression pattern in token-level alignment signal.
    \item We demonstrate that different alignment behaviors have different durability profiles: formatting instructions degrade significantly after 10--15 turns while safety refusals remain robust through 24 turns, providing empirical support for a depth-stratified view of alignment.
    \item We discover that multi-turn conversation structure amplifies instruct--base divergence relative to concatenated equivalents, suggesting that turn markers and chat templates actively modulate alignment beyond their role as formatting conventions.
\end{itemize}
